% Job posting: https://ca.linkedin.com/jobs/view/vision-software-developer-ii-at-lmi-technologies-4466393002
\documentclass[11pt]{article}
\usepackage[left=1in,top=1in,right=1in,bottom=1in]{geometry}
\usepackage[hidelinks]{hyperref}
\usepackage{setspace}
\usepackage[T1]{fontenc}
\usepackage{newtxtext,newtxmath}
\ifdefined\pdfgentounicode
  \input{glyphtounicode}
  \pdfgentounicode=1
\fi
\setstretch{1.1}
\pagenumbering{gobble}
\begin{document}
\pagestyle{empty}

\begin{flushright}
Nishal Stanislaus Silva, PhD \\
nishal.silva@hotmail.com \\
+1 613 200 2030 \\
Toronto, ON
\end{flushright}

\noindent Dear Hiring Manager,

\vspace{0.5em}
\noindent For 5.5 years at MAS Holdings, South Asia's largest apparel manufacturer, I built and
owned the computer-vision systems that decided whether a garment passed or failed on the factory
floor. I automated multilingual care-label QC with OpenCV template and feature matching and an
explainable defect overlay, iterating from a handheld scanner to an industrial camera to a
conveyor-mounted system with PLC-driven reject logic, and cut inspection time by 99.5\%. Separately,
I built a loom-side fabric-defect detector using a microscopic camera array, solving the image
registration and warping needed to reconcile overlapping fields of view, which cut operator
headcount by 50\%. LMI Technologies builds the 3D machine-vision sensors that make this class of
inspection possible at scale, and this is the exact problem space I have spent my career shipping
solutions into.

\vspace{0.5em}
\noindent A recurring thread in that work has been geometric image processing under real production
constraints. For MAS's on-demand garment-printing line, I designed the algorithm that aligns a
camera-captured garment's geometry to a print warp - a 3D-to-2D geometric registration problem - and
built the RIP integration end to end; it is now a commercial product deployed across four countries.
That is a 2D analogue of the calibration and projection problems that 3D scanning and structured-
light systems solve directly, and I bring the same discipline in C++ and Python, plus experience
deploying vision and ML systems on embedded and resource-constrained hardware, from Raspberry
Pi-class devices to production-line PLCs.

\vspace{0.5em}
\noindent I have not worked directly with structured-light or 3D-sensor hardware, and I want to be
upfront about that. What I bring instead is the surrounding discipline: taking a vision algorithm
from prototype to a reliable piece of factory-floor software, and the geometric-registration
intuition that transfers directly to 3D scanning problems. I would expect to be productive quickly
on LMI's software stack and to contribute meaningfully to inspection accuracy and robustness from
early on.

\vspace{0.5em}
\noindent I am a Canadian permanent resident and do not require sponsorship. I am based in Toronto,
open to relocating to Burnaby/Vancouver, and available immediately. I would welcome the chance to
discuss how my industrial vision background can contribute to LMI's sensor and software teams.

\vspace{1em}
\noindent Sincerely, \\
Nishal Stanislaus Silva

\end{document}
